\section{S1. Site-level carceral exposure parameters}

For the site-level sensitivity of §S2 we use the nine United States counties hosting PURPOSE 4 (NCT06101342), chosen because it is the cleanest available anchor for carceral geography in an HIV prevention trial. The registry record describes a phase 2, open-label, multicentre, randomised study of the pharmacokinetics and safety of twice-yearly subcutaneous lenacapavir for pre-exposure prophylaxis in people who inject drugs in the United States, with eligibility from 18 years and no upper bound, nine locations, and \textbf{181 participants enrolled}; it began in December 2023 and reached actual primary completion in July 2026.

Two things follow. The ≥18 frame is why adult PWID mortality from ALIVE is the appropriate source in §3.4 rather than a younger-cohort estimate. And at 181 participants across nine sites --- roughly twenty each --- \textbf{no site-level empirical claim would be supportable from this trial even in principle}, which is part of why §S2 is a break-even calculation rather than a correction. \textbf{No efficacy estimate from this trial is corrected or commented on}; it is a phase 2 pharmacokinetics and safety study and reports none.

The nine registry locations map to counties as follows, and the mapping is one-to-one: Los Angeles and San Diego, California; Miami, Florida (Miami-Dade); Baltimore, Maryland (Baltimore City); Newark, New Jersey (Essex); The Bronx, New York; Philadelphia, Pennsylvania; Houston, Texas (Harris); and Morgantown, West Virginia (Monongalia).

Jail and prison counts are taken at county level for 2019, the last year with harmonized county-level estimates across all nine counties. \textbf{2019 is a fixed pre-pandemic structural anchor and is not assumed to be conservative} --- the data disprove that reading. Post-2019 jail trajectories are heterogeneous in both magnitude and direction: relative to 2019, jail populations stand at 0.46 in the Bronx and 0.70 in San Diego, but 1.02 in Miami-Dade, 1.05 in Monongalia and 1.09 in Essex. Four of nine counties are at or above their 2019 level, so a single national multiplier would have been wrong in direction for them. A secondary analysis updates the jail component with the most recent local data while retaining the 2019 prison component, county-level post-2019 prison counts being unavailable.

Incarceration rates are published against total or 15--64 populations while trial eligibility is 18+, so a denominator conversion is required. A national 15--64 to 18+ ratio of 0.8333 was replaced with county-specific ratios from Census Population Estimates 2019, which range from 0.834 (Miami-Dade) to 0.908 (Harris). The national factor systematically understated exposure in counties with younger adult age structures, by up to 12.4\%. Consistent with §3.5, the correction is applied to the denominator only.

\begin{center}\rule{0.5\linewidth}{0.5pt}\end{center}

\section{S2. Site-specific break-even relative acquisition hazards}

To indicate the range of \(\eta\) at which the composed boundary would cross unity in real catchments, we evaluated nine United States counties with published jail and prison occupancy, using county-specific age denominators (Table S3b). Break-even \(\eta\) ranged from 0.04 (San Diego) to 0.55 (Baltimore City); in two counties, Bronx and Miami-Dade, the boundary never crosses unity for any \(\eta\in[0,1]\).

\begin{table}[htbp]
  \centering
  \small
  \caption*{\textbf{Table S3b.} Break-even $\eta$ by site county. Illustrative; not an epidemiologic claim about any county.}
  \label{tab:table_s3b}
  \begin{tabular}{lrrrrrr}
    \toprule
    site county & $q_J$ & $q_P$ & $r^\star$ at $\eta=0$ & $r^\star$ at $\eta=0.25$ & $r^\star$ at $\eta=0.5$ & break-even $\eta$ \\
    \midrule
    Bronx, NY & 0.0188 & 0.0262 & 0.9833 & 0.9607 & 0.9385 & never \\
    San Diego, CA & 0.0231 & 0.0318 & 1.0042 & 0.9760 & 0.9484 & 0.037 \\
    Miami-Dade, FL & 0.0185 & 0.0300 & 0.9857 & 0.9625 & 0.9396 & never \\
    Los Angeles, CA & 0.0206 & 0.0490 & 1.0098 & 0.9801 & 0.9510 & 0.082 \\
    Harris (Houston), TX & 0.0248 & 0.0586 & 1.0351 & 0.9986 & 0.9629 & 0.240 \\
    Monongalia, WV & 0.0294 & 0.0260 & 1.0229 & 0.9898 & 0.9573 & 0.172 \\
    Essex (Newark), NJ & 0.0318 & 0.0526 & 1.0569 & 1.0145 & 0.9731 & 0.337 \\
    Philadelphia, PA & 0.0363 & 0.0886 & 1.1124 & 1.0544 & 0.9985 & 0.493 \\
    Baltimore City, MD & 0.0385 & 0.1067 & 1.1426 & 1.0760 & 1.0121 & 0.549 \\
    \bottomrule
  \end{tabular}
\end{table}

The trial enrolled 181 participants across those nine sites, so roughly twenty each; no site-level empirical claim would be supportable from it even in principle. \textbf{This is a sensitivity range, not an epidemiologic claim about any site.} It states the value \(\eta\) would have to take for the two selection mechanisms to stop cancelling, given that county's carceral occupancy. It does not assert that \(\eta\) takes that value anywhere, and no trial estimate is corrected on its basis. Its purpose is to show that the break-even value lies inside the plausible interval implied by the only available incidence comparison, so the question is empirical rather than hypothetical.

\begin{center}\rule{0.5\linewidth}{0.5pt}\end{center}

\section{S3. Comparison with covariate-reweighting methods}

Covariate transport by reweighting --- matching the survey population to the trial-eligible population on measured covariates --- addresses a different failure and does not remove this one. The comparison uses the structure of their estimator as implemented in their own released code at commit \texttt{0581bb6} rather than as inferred from the paper: the weight is a logistic-regression predicted probability of trial membership fitted among the HIV-negative, and it multiplies numerator and denominator alike, so it cancels wherever it is constant and acts only between covariate strata. In a population where the observable and target covariate distributions coincide, which is the counterfactual-placebo case of interest, reweighting removes 0\% of the bias: the naive and reweighted estimates are identical at 0.0387 against a truth of 0.0400, both attenuated by 3.3\% (Table S2). Where the distributions differ, reweighting performs as designed, removing 91.3\% and 96.8\% of a much larger composition-driven bias in the two enriched cases.

\begin{table}[htbp]
  \centering
  \small
  \caption*{\textbf{Table S2.} Covariate reweighting against the duration mechanism, three target-population cases.}
  \label{tab:table_s2}
  \begin{tabular}{lrrrrrr}
    \toprule
    case & truth & naive & reweighted & naive bias & reweighted bias & bias removed (\%) \\
    \midrule
    same X distribution (counterfactual-placebo case) & 0.04000 & 0.03868 & 0.03868 & -0.0329 & -0.0329 & 0.0 \\
    target enriched for high-removal & 0.04800 & 0.02741 & 0.04620 & -0.4290 & -0.0375 & 91.3 \\
    target enriched for low-removal & 0.02800 & 0.04620 & 0.02741 & +0.6501 & -0.0211 & 96.8 \\
    \bottomrule
  \end{tabular}
\end{table}

The reason is structural. Reweighting corrects the \emph{composition} of the sampled population; it cannot correct the within-stratum duration component, which is below unity in every stratum, so any weighted average of within-stratum factors remains below unity. The two methods are complementary rather than alternative: reweighting for covariate imbalance, the weight \(w_t\) for eligibility dynamics.
